\documentclass[12pt]{article}

\usepackage[UTF8]{ctex}
\usepackage[margin=2.6cm]{geometry}
\usepackage{amsmath, amssymb, booktabs, float, graphicx, array, tabularx}
\usepackage{hyperref}
\usepackage{setspace}
\usepackage{titlesec}

\titleformat{\section}{\centering\heiti\zihao{4}}{\thesection}{1em}{}
\titleformat{\subsection}{\heiti\zihao{-4}}{\thesubsection}{1em}{}
\titleformat{\subsubsection}{\heiti\zihao{-4}}{\thesubsubsection}{1em}{}
\titlespacing*{\section}{0pt}{1.2ex}{1.0ex}
\titlespacing*{\subsection}{0pt}{1.0ex}{0.8ex}
\titlespacing*{\subsubsection}{0pt}{0.8ex}{0.6ex}

\makeatletter
\renewcommand{\maketitle}{
\begin{center}
{\heiti\zihao{-3}\@title\par}
\end{center}
\vspace{1em}
}
\makeatother

\title{面向智慧社区的多无人机-物业人员联合巡检优化模型}

\begin{document}

\pagestyle{plain}
\pagenumbering{arabic}
\setcounter{page}{1}
\singlespacing
\songti\zihao{-4}

\maketitle

\begin{abstract}
针对智慧社区中多无人机与物业人员联合巡检场景，本文构建一套两阶段闭环优化框架。第一阶段面向多架同构无人机，联合优化巡检目标点分配、任务批次、飞行访问顺序与累计悬停时间分配，在满足基础巡检、单次电量与运营时长约束的前提下，尽量缩短无人机巡检阶段完成时间，并兼顾总飞行能耗、高优先级点提前完成与多机负载均衡。第二阶段在空中巡检结果基础上，根据各目标点累计停留时间是否达到直接确认阈值，判定由无人机直接确认或转交物业复核，并进一步优化地面复核路径，从而最小化总闭环完成时间。基于当前附件数据的实验结果表明，$K=4$ 配置下的 ALNS 联合搜索已能将总闭环时间由 3540.7 s 进一步压缩到 2997.6 s，同时把直接确认节点由 7 个提升到 12 个、人工复核节点由 9 个降到 4 个，并将任务切分由 4 条主航线扩展为 8 条更短任务，体现出明显的空地协同降时效果。本文给出变量定义、目标函数与关键约束的统一表达，并提出后续基于精确模型与启发式算法结合的求解路线，为社区巡检资源配置与阈值设计提供定量支持。

\par\noindent{\heiti 关键词：} 多无人机巡检；联合闭环优化；直接确认；自适应大邻域搜索；空地协同
\end{abstract}

\clearpage

\section{问题重述}
题目给出社区巡检目标点、无人机飞行时间与能耗矩阵、基础悬停时长、直接确认阈值、物业地面复核时间等数据，要求建立多无人机与物业人员联合巡检优化模型。题目分为两个层次。

问题一要求在给定无人机数量 $K$ 下，确定每架无人机执行多少次任务、每次任务访问哪些目标点、访问顺序如何安排，以及各目标点的累计停留时间如何在不同无人机任务之间分配。模型需要同时满足基础空中巡检要求、单次飞行能耗上限和运营时长约束，并尽量缩短无人机阶段完成时间。

问题二在问题一基础上进一步考虑目标点是否能够由无人机直接确认。若某目标点累计停留时间达到直接确认阈值，则该点由无人机直接闭环；否则，该点需要由物业人员在无人机巡检结束后前往现场复核。因此，模型需要进一步决定哪些点由无人机直接确认、哪些点转入物业复核，并优化物业人员的地面复核顺序与通行路径，以缩短整体闭环完成时间。

从运作逻辑上看，本题的核心不是单独优化空中或地面阶段，而是在二者之间建立有效权衡：增加无人机悬停时间可以提升直接确认比例、减少人工复核，但也可能延长空中巡检阶段；相反，压缩无人机停留时间虽然有利于缩短飞行阶段，却可能增加物业复核压力。因此，本文将问题刻画为一个空地联合的多目标优化问题。

\section{模型假设与符号说明}

\subsection{基本假设}
为保证模型可解且不偏离题目本意，作如下假设：

\begin{enumerate}
	\item 所有无人机同构，飞行性能、悬停能耗与换电时间一致。
	\item 每次无人机任务均从社区服务中心出发，并在任务结束后返回服务中心。
	\item 无人机在任务执行过程中不考虑空域冲突与排队起降，仅考虑题目给出的飞行时间、飞行能耗、悬停能耗和换电时间。
	\item 目标点的巡检效果仅由累计停留时间决定，不区分由哪架无人机完成。
	\item 物业人员在无人机完成空中巡检后统一出发进行复核，不与无人机阶段并行。
	\item 物业复核点的服务时间与通行时间可加总表示，且地面复核过程不再受无人机运营时长约束。
\end{enumerate}

\subsection{主要符号}

\begin{table}[H]
\centering
\begin{tabular}{>{\raggedright\arraybackslash}p{3cm} >{\raggedright\arraybackslash}p{9.5cm}}
\toprule
符号 & 含义 \\
\midrule
$S=\{1,2,\dots,n\}$ & 巡检目标点集合 \\
$K$ & 无人机数量 \\
$T_{ij}^{u}$ & 无人机从节点 $i$ 飞行到节点 $j$ 的时间 \\
$E_{ij}^{u}$ & 无人机从节点 $i$ 飞行到节点 $j$ 的能耗 \\
$\bar{E}$ & 单次任务可用电量上限 \\
$\tau_b$ & 无人机返航后的换电时间 \\
$b_s$ & 目标点 $s$ 的基础空中巡检时间要求 \\
$h_s$ & 目标点 $s$ 的无人机直接确认累计停留阈值 \\
$T_s$ & 目标点 $s$ 的累计停留时间 \\
$p_s$ & 目标点 $s$ 的优先级权重 \\
$G_{ij}$ & 物业人员地面通行时间 \\
$T_u$ & 无人机巡检阶段完成时间 \\
$T_g$ & 物业复核阶段完成时间 \\
$T$ & 总闭环完成时间，满足 $T=T_u+T_g$ \\
$y_s$ & 若目标点 $s$ 由无人机直接确认，则 $y_s=1$，否则为 0 \\
\bottomrule
\end{tabular}
\end{table}

\section{数据理解与预处理思路}
题目配套数据表包括目标点基础信息、无人机参数、飞行时间矩阵、飞行能耗矩阵、物业复核点信息以及地面通行时间矩阵。结合题面说明，预处理将围绕以下几个方面展开。

首先，对节点编号、坐标、优先级、基础悬停时间、直接确认阈值、人工现场服务时间进行一致性校验，确保不同工作表中的点位索引可以一一对应。其次，对飞行时间矩阵与飞行能耗矩阵进行对称性、缺失值和异常值检查，识别是否存在不可达边或显著偏离平均水平的边。再次，对地面复核点数据与 GroundTime 矩阵进行核对，确保物业复核路径的建模节点与空中巡检节点之间存在稳定映射关系。

在正式求解前，还应生成若干辅助量：例如每个目标点的单位悬停收益、按优先级修正后的直接确认收益、每架无人机的有效工作窗口，以及用于衡量负载均衡的任务总时间或总服务量。上述变量既可用于构造目标函数，也可为后续启发式算法提供排序依据。

\subsection{附件字段映射}
根据对附件 \texttt{2026C数据.xlsx} 的实际读取，本文后续计算主要使用以下工作表与字段。

\begin{enumerate}
	\item \texttt{UAV\_Params}：包括 \texttt{K\_max}、\texttt{effective\_energy\_limit\_J}、\texttt{hover\_power\_J\_per\_s} 等关键参数，用于约束无人机数量、电量上限、悬停能耗和运营时长；
	\item \texttt{NodeData}：包括 \texttt{node\_id}、\texttt{priority\_level}、\texttt{priority\_weight}、\texttt{base\_hover\_time\_s}、\texttt{direct\_confirm\_time\_s}、\texttt{manual\_point\_id}、\texttt{manual\_service\_time\_s} 等字段，用于定义每个巡检点的基础需求与空地衔接关系；
	\item \texttt{FlightTime} 与 \texttt{FlightEnergy}：分别给出节点间飞行时间矩阵和飞行能耗矩阵，是构建空中巡检路径约束的核心输入；
	\item \texttt{ManualPoints} 与 \texttt{GroundTime}：分别给出物业复核点映射、服务时间以及地面通行时间矩阵，用于构造地面复核路径模型。
\end{enumerate}

在本文已完成的基线统计中，附件共包含 16 个实际巡检目标点和 1 个服务中心节点。真实字段读取说明，题面中的变量设定与附件可以一一对应，因此后续模型不仅可做理论表达，也可直接落到程序实现与实验计算。

需要说明的是，当前附件中的 \texttt{effective\_energy\_limit\_J} 为空值，因此本文在代码中按“电池总能量减安全余量”回推出有效能量上限；同时 \texttt{GroundTime} 工作表目前为空矩阵，因此地面阶段采用 \texttt{ManualPoints} 坐标、步行速度和绕行系数构造欧氏近似通行时间矩阵。该回退处理保证了模型能够继续完成联合闭环实验，也使后续结论具有明确的数据来源与可复现实验口径。

\section{问题一：多无人机巡检路径与停留时间优化模型}

\subsection{建模目标}
问题一的核心是确定空中巡检阶段的最优组织方式。本文拟将其刻画为一个多目标优化模型，主要考虑以下四类目标：

\begin{enumerate}
	\item 最小化无人机阶段完成时间 $T_u$；
	\item 最小化总飞行与悬停能耗；
	\item 提高高优先级目标点的提前完成程度；
	\item 提升多架无人机之间的任务负载均衡性。
\end{enumerate}

将多目标加权为单目标后，可写成如下形式：

\begin{equation}
\min J_1 = \omega_1 T_u + \omega_2 \sum E_{\text{total}} + \omega_3 B - \omega_4 P,
\end{equation}

其中，$B$ 表示负载不均衡惩罚项，$P$ 表示高优先级点提前完成收益，$\omega_1,\omega_2,\omega_3,\omega_4$ 为待定权重。

\subsection{关键约束}
问题一至少应满足以下约束：

为便于后续代码实现，进一步引入如下决策变量：若无人机 $k$ 的第 $r$ 次任务经过弧 $(i,j)$，则令二元变量 $x_{ij}^{kr}=1$，否则为 0；若无人机 $k$ 在第 $r$ 次任务中对目标点 $s$ 的悬停时间为 $z_s^{kr}$，则有

\begin{equation}
T_s = \sum_k \sum_r z_s^{kr}, \quad z_s^{kr} \ge 0.
\end{equation}

各任务的路径流守恒关系可统一表示为

\begin{equation}
\sum_j x_{ij}^{kr} - \sum_j x_{ji}^{kr} = 0,
\end{equation}

其中除服务中心外，任一被访问节点在同一任务中的进入次数与离开次数相等；服务中心则满足起点与终点闭环约束。

\begin{equation}
T_s \ge b_s, \quad \forall s \in S,
\end{equation}

即所有目标点均须完成基础空中巡检。

对于任一无人机单次任务，其飞行与悬停总能耗不得超过单次可用电量上限：

\begin{equation}
E_{\text{route}} + E_{\text{hover}} \le \bar{E}.
\end{equation}

若记单位悬停能耗率为 $\rho_h$，则有

\begin{equation}
E_{\text{hover}} = \rho_h \sum_{s \in S} z_s^{kr}.
\end{equation}

对任一无人机，其任务总持续时间由飞行、悬停和换电共同构成。若无人机 $k$ 共执行 $R_k$ 次任务，则可写为

\begin{equation}
T_u^{(k)} = \sum_{r=1}^{R_k} \left(T_{\text{route}}^{kr} + T_{\text{hover}}^{kr}\right) + (R_k-1)\tau_b.
\end{equation}

所有无人机执行完全部任务所需时间应落在运营时段之内，且系统阶段完成时间定义为各无人机完成时间的最大值：

\begin{equation}
T_u = \max_{k=1,2,\dots,K} T_u^{(k)}.
\end{equation}

另外，每次任务均需满足从服务中心出发并返回服务中心的闭环约束；一个目标点可以在多次任务中被重复访问，其累计停留时间等于所有访问中悬停时间之和。

\subsection{求解思路}
从结构上看，问题一属于带悬停时间分配与电量约束的多旅行商/车辆路径扩展问题。本文计划采用两层求解思路：对小规模实例建立混合整数规划模型，以获得基准解或近最优解；对中大规模实例采用“点位分组 + 路径构造 + 局部搜索修正”的启发式流程，以兼顾求解效率与结果稳定性。

更具体地，启发式求解将按以下步骤进行：

\begin{enumerate}
	\item 根据空间位置、优先级与基础巡检时间对目标点进行预分组，先用 K-means 生成初始任务簇，并为每个簇选取最接近质心的代表点作为无人机起始候选；
	\item 在每个任务簇内构造闭环飞行路径，并检查单次任务电量约束；
	\item 在满足基础巡检时间的前提下，为高优先级点和潜在直接确认点追加悬停时间；
	\item 通过插入、交换、跨任务重分配与 2-opt 等局部搜索算子迭代改进；
	\item 以系统完成时间、总能耗和负载均衡指标为准则筛选最终方案。
\end{enumerate}

结合当前已实现的 Python 原型，本文首先落地一版“构造式启发式”作为问题一的基线求解器。该算法不直接追求联合优化直接确认收益，而是优先保证所有目标点满足基础巡检时间要求，并在此基础上尽量压缩多机并行后的最大完工时长。其核心步骤如下：

\begin{enumerate}
	\item 先对全部目标点执行 K-means 聚类，并将每个簇中最接近质心的点作为多无人机的优先起点候选，再结合优先级权重、基础悬停时间和直接确认阈值生成多起点访问顺序；
	\item 对每个待分配目标点，同时评估两类动作：插入到某架无人机已有航线的任意位置，或为某架无人机新开一条从服务中心出发并回收的独立航线；
	\item 对每个候选动作显式计算新增飞行时间、返航时间、悬停时间与新增能耗，只保留满足单次任务能耗上限和单机运营时段上限的可行动作；
	\item 以“新增后的系统最大完工时长最小”为第一准则，以“新增后的该机总时长”和“任务总能耗”作为次级准则进行贪心选择；
	\item 当所有目标点均被分配后，输出每架无人机的任务条数、各条航线访问顺序、每个目标点的悬停时间以及由此形成的系统完工时间、总能耗和负载离散度。
\end{enumerate}

这一定义方式与题目所要求的四个核心决策完全对应：哪些无人机负责哪些点、每架无人机执行多少条任务、每条任务的访问顺序，以及各目标点基础悬停时间如何分配到对应任务中。由于本题数据规模仅为 16 个巡检点，该构造式启发式能够在很短时间内输出稳定可行解，适合作为后续局部搜索和问题二联动优化的出发点。

\section{问题二：多无人机与物业人员联合闭环优化模型}

\subsection{直接确认判定}
在问题一的基础上，引入二元变量 $y_s$ 表示目标点是否由无人机直接确认：

\begin{equation}
y_s =
\begin{cases}
1, & T_s \ge h_s, \\
0, & T_s < h_s.
\end{cases}
\end{equation}

若 $y_s=0$，则该点必须进入物业复核集合。此时，空中巡检与地面复核之间形成资源替代关系：提升 $T_s$ 有助于增加直接确认比例，但也会推高空中阶段耗时；降低 $T_s$ 则可能导致更多目标点进入物业复核，从而抬高 $T_g$。

\subsection{联合优化目标}
问题二的整体目标是最小化总闭环完成时间，并兼顾高优先级点尽量由无人机直接确认、人工复核点数量减少等要求。可写为：

\begin{equation}
\min J_2 = \lambda_1 (T_u + T_g) + \lambda_2 \sum_{s \in S} p_s (1-y_s) + \lambda_3 N_g,
\end{equation}

其中 $N_g$ 表示需要人工复核的目标点数量，$\lambda_1,\lambda_2,\lambda_3$ 为权重参数。

\subsection{地面复核路径}
对于所有需要人工复核的点，物业人员从服务中心统一出发并完成一条地面复核路径。若记复核点集合为 $S_g=\{s \in S : y_s=0\}$，则问题可进一步转化为一类旅行商路径问题，其成本由地面通行时间与现场复核服务时间共同构成。物业阶段完成时间记为 $T_g$，并与空中阶段串联构成总闭环时间。

若记物业人员是否从复核点 $i$ 前往复核点 $j$ 的二元变量为 $g_{ij}$，点 $s$ 的现场复核服务时间为 $\tau_s^g$，则地面阶段目标可写为

\begin{equation}
T_g = \sum_{i}\sum_{j} G_{ij} g_{ij} + \sum_{s \in S_g} \tau_s^g.
\end{equation}

因此，问题二的关键在于同时决定“空中多停留一点”和“地面少复核一点”之间的边际替代关系。若某些高优先级点在地面复核中通行代价较高，则更应在空中阶段优先满足其直接确认阈值。

\section{算法框架与实验设计}
为保证模型既能在论文中给出清晰数学表达，也能在竞赛时间内落地实现，本文拟采用如下实验框架：

\begin{enumerate}
	\item 以问题一为基线模型，比较不同无人机数量 $K$ 下的完成时间、总能耗和负载均衡度；
	\item 在问题二中比较“仅完成基础巡检”和“兼顾直接确认能力”的两类策略；
	\item 对关键参数进行灵敏度分析，包括无人机数量 $K$、换电时间 $\tau_b$ 和直接确认阈值 $h_s$；
	\item 基于边际收益递减规律，给出推荐的无人机规模与阈值设定建议。
\end{enumerate}

从实现角度看，后续代码部分将数据读取、空中路径求解、停留时间分配、直接确认判定与地面复核路径优化封装为独立模块，以便批量运行参数实验并统一生成图表。

\subsection{实验输出设计}
为了让论文结果部分具有可直接展示的结构，本文预设如下核心输出：

\begin{enumerate}
	\item 不同无人机数量下的系统总闭环完成时间、空中阶段完成时间、地面阶段完成时间；
	\item 不同策略下的总飞行能耗、平均任务时长与最大任务时长；
	\item 无人机直接确认比例、人工复核点数量以及高优先级点直接确认比例；
	\item 对 $K$、$\tau_b$ 与 $h_s$ 变化的灵敏度曲线；
	\item 推荐配置方案及其相对基线方案的改进幅度。
\end{enumerate}

建议在正式实验中设置两类对照：一类是“仅满足基础巡检时间”的保守方案，另一类是“兼顾直接确认收益”的联合优化方案。通过这两类方案的对比，可以直观体现空地协同建模的必要性。

\subsection{基线统计结论}
基于当前已完成的数据脚手架，本文先对附件数据做一轮基线统计。结果显示，附件共包含 16 个实际巡检目标点，基础悬停时间均值为 49.375 秒，而直接确认阈值均值达到 325.625 秒，二者差距显著。在“仅满足基础巡检时间”的基线方案下，16 个目标点中没有任何一个点能够直接达到无人机确认阈值，直接确认比例为 0。这说明若不在问题一的基础上进一步优化悬停时间分配，问题二中的物业复核阶段将不可避免地承担几乎全部闭环确认任务。

与此同时，基线统计还表明，无论是完成基础巡检还是完成直接确认，对 16 个目标点逐点考察时其单点能耗均未超过附件给定的有效能量上限 135000J。这意味着当前问题的主要矛盾并非“单点是否能飞到”，而是“在有限运营时长和多目标点竞争下，如何分配总悬停资源以尽量减少人工复核”。因此，附件中给出的单点往返飞行能耗字段以及直接确认总能耗字段，不仅能够用于快速判定单点可行性，也能为后续启发式算法中的优先级排序和初始解构造提供依据。

\subsection{问题一最终采用的最优算法与结果展示}
在补全缺失飞行边并引入 K-means 起点候选之后，本文对问题一两类可行构造算法做了同口径比较：一类是以最小化系统最大完工时长为主的构造式贪心方案，另一类是带层次滚动分组的重构方案。最终结果表明，当前最优算法并非对所有 $K$ 都相同，而是呈现出“低到中等 $K$ 下构造式贪心更优、$K=4$ 时层次滚动重构更优”的分段特征。实验结果见表 \ref{tab:problem1-heuristic-result}。

\begin{table}[H]
\centering
\caption{问题一最终采用的最优算法在不同无人机数量下的结果对比}
\label{tab:problem1-heuristic-result}
\begin{tabular}{cccccc}
\toprule
$K$ & 航线条数 & 最大完工时长/s & 最大完工时长/min & 总能耗/J & 负载标准差/s \\
\midrule
1 & 2 & 1312.2 & 21.87 & 210813 & 0.00 \\
2 & 2 & 537.4 & 8.96 & 219924 & 8.94 \\
3 & 3 & 361.3 & 6.02 & 220008 & 6.84 \\
4 & 4 & 303.4 & 5.06 & 239366 & 14.87 \\
\bottomrule
\end{tabular}
\end{table}

由表 \ref{tab:problem1-heuristic-result} 可见，随着无人机数量从 1 增加到 4，系统最大完工时长从 1312.2 s 快速下降到 303.4 s，表明多机并行对缩短空中巡检阶段依然具有显著效果。其中，从 1 架增至 2 架的改进最明显，完工时间下降约 59.0\%；继续从 2 架增至 3 架和 4 架时，仍有改善，但边际收益开始递减。更重要的是，$K=4$ 的最优结果已经不再来自早期构造式贪心，而是来自层次滚动重构方案，这说明当无人机数量提高到 4 架时，仅靠局部贪心插入已不足以充分利用更细粒度的并行空间，需要显式引入分组重构与滚动调整机制。

另一方面，总悬停时间在各组实验中保持不变，均为 790 s，这是因为当前问题一阶段仍以满足基础巡检需求为主，尚未在空中路径构造阶段直接追求更多直接确认。因此，不同 $K$ 的能耗差异主要来自路径结构变化而非服务强度变化。结果显示，总能耗处于 21.1 万 J 到 23.9 万 J 区间内，说明在本题数据规模下，问题一的主要权衡仍是“更短的最大完工时长”与“更均衡的航线切分”，而不是能耗约束本身成为主导瓶颈。

负载均衡方面，$K=2$ 时负载标准差约为 8.94 s，说明两机分工已经较为均衡；$K=3$ 时标准差进一步降到 6.84 s，说明三机配置在当前 16 点规模下能同时兼顾完工时间和负载均衡；当 $K$ 增至 4 时，最终采用的层次滚动方案将标准差控制在 14.87 s，相比早期 23.09 s 的构造式结果明显改善，说明更强的路径重构机制确实能够缓解“四机并行但个别无人机过短”的失衡现象。

从航线结构看，当前最终采用的问题一算法已能自动将 16 个目标点压缩为少量可行航线。例如，$K=1$ 时形成 2 条可行闭环航线；$K=2$ 时，两架无人机各承担 1 条长度接近的主航线，系统最大完工时长降至 537.4 s；$K=3$ 时则形成 3 条持续时间接近的并行航线；$K=4$ 时进一步形成 4 条更短的任务链，并通过层次滚动重构获得当前问题一的最优空中阶段表现。所有方案均满足单次任务能耗上限 135000 J 以及单机运营时段 2600 s 的约束，因此该组合式策略已具备作为问题一正式实验基线的可用性。

为更直观展示不同 $K$ 下的任务空间分布，图 \ref{fig:problem1-route-panels} 给出了问题一最终航线的二维投影。可以看到，随着无人机数量增加，原本需要单机跨越社区南北两端的大范围闭环被逐步拆解为若干更短的局部链条，这与表 \ref{tab:problem1-heuristic-result} 中“最大完工时长快速下降、但边际收益逐渐减弱”的结论保持一致。
其中，$K=1$ 面板中出现两条折线，并不意味着有两架无人机，而是表示同一架无人机在该配置下需要先后执行两次完整任务；图中的每一条线对应一条完整 route，而不是一架无人机。

\begin{figure}[H]
\centering
\includegraphics[width=0.95\linewidth]{outputs/figures/c_problem_problem1_route_panels.png}
\caption{问题一在不同无人机数量下的代表航线路径总览}
\label{fig:problem1-route-panels}
\end{figure}

\subsection{K-means 起点选取消融实验}
为了验证 K-means 起点候选是否真的带来收益，本文将当前方案与上一版未使用 K-means 起点的缓存航线基线做同口径对比，结果见表 \ref{tab:problem1-kmeans-ablation}，可视化见图 \ref{fig:problem1-kmeans-ablation}。

\begin{table}[H]
\centering
\caption{K-means 起点选取前后问题一结果对比}
\label{tab:problem1-kmeans-ablation}
\begin{tabular}{ccccc}
\toprule
$K$ & K-means 前完工时长/s & K-means 后完工时长/s & 绝对改进/s & 相对改进 \\
\midrule
1 & 1491.0 & 1312.2 & 178.8 & 12.0\% \\
2 & 582.2 & 537.4 & 44.8 & 7.7\% \\
3 & 406.7 & 361.3 & 45.4 & 11.2\% \\
4 & 322.9 & 322.4 & 0.5 & 0.1\% \\
\bottomrule
\end{tabular}
\end{table}

\begin{figure}[H]
\centering
\includegraphics[width=0.92\linewidth]{outputs/figures/c_problem_problem1_kmeans_ablation.png}
\caption{K-means 起点选取前后问题一最大完工时长对比}
\label{fig:problem1-kmeans-ablation}
\end{figure}

消融实验表明，K-means 起点策略在 $K=1,2,3$ 下都带来稳定改进，最大完工时长分别下降 178.8 s、44.8 s 和 45.4 s，降幅约为 12.0\%、7.7\% 和 11.2\%。更重要的是，四组实验中的航线条数保持不变，因此这部分收益并不是通过增加任务条数取得的，而是来源于更合理的初始簇划分和更短的任务内部访问顺序。对 $K=4$，K-means 仍带来 0.5 s 的微小改进，但收益已明显趋于饱和，说明在当前 16 个点规模下，起点优化的边际收益主要集中在 $K=1$ 到 $K=3$。

结合聚类结果可见，$K=3$ 时三个簇的代表起点分别落在节点 12、14 和 5 附近，能够较好覆盖南侧、北西侧和东侧建筑群。这样做的直接效果是降低了多架无人机在首次出发阶段的重复绕行距离，也为后续问题二的状态引导追加悬停留下了更均衡的空中余量。因此，本文后续关于问题二的实验全部采用这一版 K-means 起点构造结果作为新的问题一输入。

\subsection{问题二首轮闭环对比结果}
在问题一可行航线的基础上，本文进一步实现了问题二的联合闭环求解链路，并比较三类策略：\textbf{baseline} 表示仅满足基础巡检时间；\textbf{optimize\_hover} 表示基于边际收益的追加悬停贪心；\textbf{ga\_hover} 表示在候选点集合上执行遗传算法全局搜索。三类方案在不同无人机数量下的总闭环完成时间见表 \ref{tab:problem2-round1}。

\begin{table}[H]
\centering
\caption{问题二首轮闭环算法对比结果}
\label{tab:problem2-round1}
\begin{tabular}{ccccc}
\toprule
$K$ & baseline/s & 边际收益贪心/s & GA/s & 最优方案直接确认数 \\
\midrule
1 & 5706.8 & 5642.4 & 5642.4 & 1 \\
2 & 4798.0 & 4717.4 & 4717.4 & 1 \\
3 & 4622.5 & 4183.8 & 4112.2 & 4 \\
4 & 4538.7 & 3796.2 & 3693.1 & 6 \\
\bottomrule
\end{tabular}
\end{table}

表 \ref{tab:problem2-round1} 表明，问题二的主要收益已不再来自单纯压缩空中阶段，而是来自“增加直接确认比例、减少物业复核点数”所带来的闭环缩短。对 $K=1,2$，GA 与边际收益贪心结果基本一致，说明此时可追加悬停余量有限，局部选择已足以逼近较优方案；但当 $K$ 提高到 3 或 4 时，GA 分别将总闭环完成时间进一步压缩到 4112.2 s 和 3693.1 s，较当前贪心再降低 71.6 s 和 103.1 s。这说明当多机并行已经释放出一定空中余量后，全局搜索比单步贪心更容易找到“多停一点但少走很多人工路”的组合。

从直接确认数量看，当前最优策略在 $K=3$ 时可直接确认 4 个目标点，在 $K=4$ 时可直接确认 6 个目标点，对应人工复核点数分别降为 12 和 10。与问题一相比，问题二已经把评价重点从“空中最短完工时间”推进到了“空地联合闭环最短完成时间”，并验证了直接确认阈值确实是影响物业复核负担的关键变量。

\subsection{问题二第二轮状态引导架构}
在首轮闭环实验基础上，本文进一步将问题二的追加悬停决策统一改写为状态引导函数 $F(u,v,E)$。其中，$u$ 表示当前航线上的已访问点，$v$ 表示下一候选追加悬停点，$E$ 表示无人机到达 $u$ 时的剩余可用能量。若从当前状态转向 $v$ 后会破坏返航充电或运营时长约束，则该转移直接视为无效；否则，用当前附件数据中可直接计算的六类量构造状态评分：地面节省收益、点位优先级、剩余能量松弛、剩余时长松弛、航线推进位置和悬停缺口惩罚。具体实现为

\begin{equation}
\begin{aligned}
F(u,v,E) &= 0, && \text{若转移后不可返航；} \\
F(u,v,E) &= w_1 \tilde{G}(u,v)+w_2 \tilde{P}(v)+w_3 \tilde{S}_E(v)
		   + w_4 \tilde{S}_T(v) \\
		  &\quad + w_5 \tilde{R}(v)-w_6 \tilde{H}(v), && \text{否则。}
\end{aligned}
\end{equation}

其中，$\tilde{G}(u,v)$ 表示追加确认点 $v$ 后所节省的地面服务时间与路径绕行时间，$\tilde{P}(v)$ 表示点位优先级，$\tilde{S}_E(v)$ 与 $\tilde{S}_T(v)$ 分别表示追加后剩余能量和剩余运营时长的相对松弛度，$\tilde{R}(v)$ 表示该点在既有航线中的推进位置，$\tilde{H}(v)$ 表示从当前悬停时长补足到直接确认阈值所需的额外悬停缺口。每次一旦选择某个点进行追加悬停，系统状态便从 $(u,E)$ 转移到新的 $(v,E')$，并按新的剩余能量与剩余时长重新计算全部候选点的 $F$ 值。这样既保留了你提出的“状态递推式指引”思想，又避免了在竞赛时间内离散枚举完整三维状态表的额外复杂度。

在当前实现中，默认策略 \textbf{f\_default} 使用固定权重向量直接驱动上述评分，而 GA、PSO、ACO 三类元启发式则不再直接搜索“选哪些点”，而是统一搜索这组权重参数，使三类算法真正建立在同一个 $F(u,v,E)$ 架构之上。

\subsection{问题二第三轮默认 F 策略实验}
在引入统一状态引导函数之后，首先比较基础闭环方案与默认 F 策略 \textbf{f\_default}。实验仍聚焦在问题一已经显示出较强工程价值的 $K=3,4$ 配置下，结果见表 \ref{tab:problem2-round2}。

\begin{table}[H]
\centering
\caption{问题二默认 F 策略实验结果}
\label{tab:problem2-round2}
\begin{tabular}{ccccc}
\toprule
 $K$ & baseline/s & 默认 F/s & 直接确认数 & 相对 baseline 降幅 \\
\midrule
3 & 4606.7 & 4180.8 & 4 & 9.2\% \\
4 & 4511.9 & 3534.5 & 7 & 21.7\% \\
\bottomrule
\end{tabular}
\end{table}

表 \ref{tab:problem2-round2} 表明，在更新问题一路线后，单纯将问题二写成“状态引导函数驱动”依然能稳定获得明显收益。对 $K=3$，默认 F 策略可将总闭环时间从 4606.7 s 降至 4180.8 s，直接确认 4 个点；对 $K=4$，则可进一步降至 3534.5 s，并直接确认 7 个点。与问题一相比，这一轮实验说明：问题二的核心不再是单纯压缩空中飞行时间，而是根据“当前航线位置 + 剩余能量 + 地面节省收益”判断某个点是否值得继续空中确认。与此同时，K-means 起点构造缩短了空中路径，也放大了高 $K$ 场景下状态引导追加悬停的收益空间。

\subsection{问题二第四轮 F 函数三算法学习实验}
在默认 F 策略的基础上，本文继续使用 GA、PSO、ACO 三类元启发式学习同一组 $F(u,v,E)$ 权重。这样三类算法的差异只体现在“如何更新 F 的权重”，而不再混杂不同的选点编码方式。实验结果见表 \ref{tab:problem2-round3}。

\begin{table}[H]
\centering
\caption{问题二 F 函数三算法学习实验结果}
\label{tab:problem2-round3}
\begin{tabular}{cccccc}
\toprule
$K$ & 默认 F/s & F-GA/s & F-PSO/s & F-ACO/s & 最优策略 \\
\midrule
3 & 4180.8 & 3982.5 & 3982.5 & 3946.6 & ACO \\
4 & 3534.5 & 3534.5 & 3534.5 & 3534.5 & 并列 \\
\bottomrule
\end{tabular}
\end{table}

表 \ref{tab:problem2-round3} 给出两个新的现象。第一，对 $K=3$，F-GA 与 F-PSO 都能把闭环时间进一步降到 3982.5 s，而 F-ACO 还能继续压低到 3946.6 s，并将直接确认数提升到 5 个。这说明在当前数据规模下，$K=3$ 仍然存在可以被学习到的更优 F 权重结构，而 ACO 在该结构上最具优势。第二，对 $K=4$，默认 F 与三类元启发式全部收敛到同一结果 3534.5 s，说明在新的问题一路线下，$K=4$ 对应的追加悬停决策已经非常稳定，进一步学习 F 权重并不会带来额外收益。

从最终效果看，$K=3$ 更适合作为“有限资源下追求更优策略学习”的配置，因为 F 权重学习仍能带来 5\% 以上的额外闭环压缩；而 $K=4$ 则更像是“更高资源投入下的稳定饱和配置”，默认 F 已足够接近最优。由此可见，在当前 16 点规模和既定数据结构下，F 函数学习型 ACO 主要在 $K=3$ 场景最值得继续深化。

在此基础上，本文又补做了一轮“问题一路径重排 + 问题二 F 引导”的联合实验。具体做法是：在不改变问题一无人机分组与航线条数的前提下，尝试在每条新航线内部按照当前默认 F 权重重新排序，再重新执行问题二的 F 引导追加悬停决策，并将该联合方案在 $K=1$ 至 $K=4$ 上与固定路线 F 做并列比较。实验结果显示，联合重排并非总是有效：对 $K=3$，联合重排将总闭环时间从 4180.8 s 进一步压缩到 4048.6 s，额外改善 132.2 s；但对 $K=2$ 和 $K=4$，联合重排分别轻微变差 38.9 s 和 0.5 s；对 $K=1$，则由于单机长航线结构受限并伴随回退，无额外收益。这个结果说明：K-means 起点和 F 感知构造已经显著改写了问题一对问题二的传导关系，问题一路线并非“对问题二没有影响”，但其改进效果具有明显的 $K$ 依赖性，其中 $K=3$ 是当前最值得继续深化的联合优化配置。

\subsection{评价指标}
为了保证与前文实验结果保持可比，本文保留旧评价标准，同时引入一套更强调“点位优先级覆盖”和“人工复核负担”的补充评价标准。

旧评价标准仍包括总闭环完成时间 $T=T_u+T_g$、直接确认比例和负载均衡度。其中，

\begin{equation}
\text{DirectRate} = \frac{\sum_{s \in S} y_s}{n},
\end{equation}

\begin{equation}
\text{Balance} = \sqrt{\frac{1}{K} \sum_{k=1}^{K} \left(T_u^{(k)} - \frac{1}{K}\sum_{k=1}^{K} T_u^{(k)}\right)^2 }.
\end{equation}

其中，DirectRate 用于衡量无人机直接闭环能力，Balance 用于衡量多机负载离散程度，数值越小表示任务分配越均衡。

在此基础上，本文新增三项补充指标。第一项是优先级加权直接确认率

\begin{equation}
\text{PriorityDirectRate} = \frac{\sum_{s \in S} p_s y_s}{\sum_{s \in S} p_s},
\end{equation}

它衡量被无人机直接闭环的点位中，优先级权重被覆盖的比例。第二项是人工复核服务时间占比

\begin{equation}
\text{ManualServiceRatio} = \frac{\sum_{s \in S} \tau_s^g (1-y_s)}{\sum_{s \in S} \tau_s^g},
\end{equation}

用于衡量剩余人工现场服务负担。第三项是加权人工复核负担占比

\begin{equation}
\text{WeightedManualBurden} = \frac{\sum_{s \in S} p_s\tau_s^g(1-y_s)}{\sum_{s \in S} p_s\tau_s^g},
\end{equation}

它同时刻画“未被直接确认的点是否优先级高”和“这些点的人工复核时间是否长”。为了给当前若干方案做统一补充比较，定义扩展评价分数

\begin{equation}
\text{Score}_{+} = 0.55\frac{T}{T^{\text{base}}_K} + 0.25\left(1-\text{PriorityDirectRate}\right) + 0.20\,\text{ManualServiceRatio},
\end{equation}

其中 $T^{\text{base}}_K$ 表示同一无人机数量 $K$ 下 baseline 方案的总闭环完成时间。这样处理后，旧标准仍负责回答“谁更快”，新标准则进一步回答“快的同时是否覆盖了更重要的点，以及是否真正降低了物业复核负担”。

\subsection{问题二在新旧标准下的补充评价}
基于当前一致口径的轻量/当前结果表，本文对 baseline、默认 F、联合重排 F 和任务次数联合方案做补充评价，结果见表 \ref{tab:problem2-expanded-eval}。需要说明的是，当前轻量 F-GA、F-PSO、F-ACO 与默认 F 在 $K=3,4$ 下收敛到相同直接确认点集，因此在扩展评价下也与默认 F 并列，不再重复列入表中。

\begin{table}[H]
\centering
\caption{问题二代表方案在新旧评价标准下的补充对比}
\label{tab:problem2-expanded-eval}
\small
\setlength{\tabcolsep}{3.5pt}
\begin{tabular}{ccccccc}
\toprule
$K$ & 方案 & \shortstack{旧标准\\闭环时间/s} & \shortstack{Priority\\DirectRate} & \shortstack{Manual\\ServiceRatio} & \shortstack{Weighted\\ManualBurden} & $\text{Score}_{+}$ \\
\midrule
3 & baseline & 4572.0 & 0.000 & 1.000 & 1.000 & 1.000 \\
3 & 默认 F & 3870.2 & 0.222 & 0.719 & 0.796 & 0.804 \\
3 & 联合重排 F & 3880.0 & 0.222 & 0.719 & 0.796 & 0.805 \\
3 & 任务次数联合 & 3870.2 & 0.222 & 0.719 & 0.796 & 0.804 \\
4 & baseline & 4533.2 & 0.000 & 1.000 & 1.000 & 1.000 \\
4 & 默认 F & 3540.7 & 0.333 & 0.607 & 0.699 & 0.718 \\
4 & 联合重排 F & 3550.1 & 0.333 & 0.607 & 0.699 & 0.719 \\
4 & 任务次数联合 & 3540.7 & 0.333 & 0.607 & 0.699 & 0.718 \\
\bottomrule
\end{tabular}
\normalsize
\end{table}

表 \ref{tab:problem2-expanded-eval} 给出两点新信息。第一，在当前结果口径下，新标准并没有推翻旧标准的排序：对 $K=3$，默认 F 与任务次数联合并列最优，联合重排 F 略差；对 $K=4$，默认 F 与任务次数联合并列最优，联合重排 F 依然略差。这说明此前基于总闭环完成时间得到的主要结论仍然稳定。第二，新标准揭示了旧标准中不够显性的短板。虽然 $K=4$ 相比 $K=3$ 在总闭环时间上进一步下降，但其 PriorityDirectRate 也仅为 0.333，ManualServiceRatio 仍有 0.607，说明当前方案的收益主要来自缩短总时间和减少部分人工复核，而不是大量消化高优先级且人工服务代价大的点位。

进一步检查当前直接确认点集可知，$K=3$ 下代表方案直接确认的点为 1、3、10、11、12，$K=4$ 下为 1、3、6、7、10、11、12。它们大多属于直接确认阈值较低或中等优先级点，因此当前所有代表方案的“高优先级点直接确认率”仍为 0。这说明把点位优先级和人工复核时间显式加入评价标准后，模型虽然仍会推荐 $K=4$，但推荐理由已不再只是“更快”，而是“在当前算法能力下它给出最好旧标准结果，同时在新标准下也没有被更复杂方案超越”。

图 \ref{fig:problem2-scheme-compare} 则把代表方案与当前已找到的 ALNS 最优结果放到同一坐标系中。该图一方面直观重现了 $K=3$ 下各代表方案几乎重合、$K=4$ 下默认 F 与任务次数联合仍构成稳定基线的事实；另一方面也可看出 ALNS 的最佳结果虽然显著低于当前 $K=4$ 基线，但这一优势还需要结合后续稳定性分析来判断是否可作为正式结论。

\begin{figure}[H]
\centering
\includegraphics[width=0.95\linewidth]{outputs/figures/c_problem_problem2_scheme_compare.png}
\caption{问题二代表方案与 ALNS 最优结果的闭环时间对比}
\label{fig:problem2-scheme-compare}
\end{figure}

\subsection{ALNS 稳定性与跨 $K$ 航线合并补充实验}
在上述结果基础上，本文进一步补做两类补充实验。第一类实验考察 ALNS 联合搜索器的稳定性：保留当前默认 F 作为评估内核，只改变 destroy-repair 搜索过程，在不同无人机数量下做轻量多种子验证。第二类实验考察一种跨 $K$ 的问题一启发：先取 $K=4$ 的 4 条短航线，将其两两配对后交给 2 架无人机顺序执行；再取 $K=2$ 的两条航线，交给 1 架无人机顺序执行，以此验证“高 $K$ 短航线方案是否能反向启发低 $K$ 方案”。

ALNS 稳定性结果见表 \ref{tab:problem2-alns-stability}。与前一版轻量验证相比，当前多种子结果已经更清晰地表明：$K=4$ 是 ALNS 最值得继续深化的配置。3 组种子全部优于统一 F 基线，最佳结果可将总闭环时间从 3540.7 s 压到 2997.6 s，最佳增益达到 543.1 s，平均增益也达到 328.6 s；相比之下，$K=3$ 虽然也出现了 3752.7 s 的更优结构，但改进只在 1 组种子上出现，$K=1,2$ 的收益则仍然较小。这意味着 ALNS 的作用已经从“证明更优结构存在”进一步推进到“在 $K=4$ 上具备较强复现性”。进一步地，本文对当前批次中表现最好的 seed=37 做了 12 次迭代的 deeper 复验，得到的最优闭环时间为 3077.4 s，对应 8 条任务、11 个直接确认点和 563.5 s 的求解耗时。该结果没有超过历史最优 2997.6 s，但复现了“通过更大邻域重构可显著优于联合基线”的总体结论，因此后续更值得继续加强的是算子权重自适应与多 seed 长跑稳定性，而不是回退到更窄的局部搜索邻域。

\begin{table}[H]
\centering
\caption{ALNS 轻量多种子稳定性补充结果}
\label{tab:problem2-alns-stability}
\begin{tabular}{cccccc}
\toprule
$K$ & 种子数 & 最佳闭环时间/s & 最佳增益/s & 平均增益/s & 改进种子比例 \\
\midrule
1 & 3 & 5291.9 & 32.5 & 10.8 & 0.33 \\
2 & 3 & 4545.2 & 5.4 & 1.8 & 0.33 \\
3 & 3 & 3752.7 & 117.5 & 39.2 & 0.33 \\
4 & 3 & 2997.6 & 543.1 & 328.6 & 1.00 \\
\bottomrule
\end{tabular}
\end{table}

跨 $K$ 航线合并实验结果见表 \ref{tab:problem1-cross-k-merge}。结果表明，这种“把高 $K$ 下的多条短航线直接合并给更少无人机顺序执行”的启发在当前数据下并不成立。对 $K=4 \rightarrow 2$，3 种两两配对后的候选最大完工时间分别为 903.7 s、883.6 s 和 895.1 s，均显著劣于当前问题一直接求得的 $K=2$ 基线 537.4 s；对 $K=2 \rightarrow 1$，顺序合并后的单机总时长为 1356.9 s，也劣于当前 $K=1$ 基线 1312.2 s。其原因在于：当前问题一的直接求解器会在更少无人机条件下同步重构分组和访问顺序，而简单地继承高 $K$ 短航线再串行化，只会额外引入换电与返航开销，无法保留高 $K$ 方案的并行优势。

\begin{table}[H]
\centering
\caption{跨 $K$ 航线合并启发补充实验结果}
\label{tab:problem1-cross-k-merge}
\small
\setlength{\tabcolsep}{3pt}
\begin{tabular}{ccccc}
\toprule
目标配置 & 候选方案 & \shortstack{启发式\\完工时间/s} & \shortstack{当前基线\\完工时间/s} & \shortstack{相对基线\\改进/s} \\
\midrule
$K=4 \rightarrow 2$ & $(1,2)|(3,4)$ & 903.7 & 537.4 & -366.3 \\
$K=4 \rightarrow 2$ & $(1,3)|(2,4)$ & 883.6 & 537.4 & -346.2 \\
$K=4 \rightarrow 2$ & $(1,4)|(2,3)$ & 895.1 & 537.4 & -357.7 \\
$K=2 \rightarrow 1$ & $(1,2)$ 串行执行 & 1356.9 & 1312.2 & -44.7 \\
\bottomrule
\end{tabular}
\normalsize
\end{table}

这两组补充实验共同说明：真正有继续深化价值的，不是简单的跨 $K$ 串并启发，而是围绕 $K=4$ 的 destroy-repair 结构继续做联合搜索设计。当前结果已经表明，问题一路由构造给出的 4 条主航线并不是闭环求解的终点，后续通过更细粒度任务切分、跨航线节点重配和更激进的直接确认推进，仍能明显继续压缩总闭环时间。

图 \ref{fig:problem2-alns-merge-visuals} 进一步把这两组补充实验转成图形展示。左图显示，当前 ALNS 的显著收益主要集中在 $K=4$，而跨 $K$ 航线合并启发的所有候选方案依旧劣于直接针对目标 $K$ 重构得到的基线方案。因此，后续优化主线应继续围绕“联合搜索如何稳定得到更优闭环解”展开，而不是回到简单的串并合并启发。

进一步地，图 \ref{fig:problem2-k4-default-vs-alns} 直接对比了 $K=4$ 下默认 F 路径与当前最优 ALNS 路径。可以看到，默认 F 仍保持 4 条较长任务，而 ALNS 最优解将其重构为 8 条更短任务，并把更多节点从“人工复核”推进为“空中直接确认”。闭环时间由 3540.7 s 进一步压缩到 2997.6 s，说明当前最优解并不是在原有路径上做小修小补，而是通过任务粒度重构与确认节点重分配同时获得收益。

若进一步把节点状态和历程看板一起纳入观察，则这种改进会更直观。图 \ref{fig:problem2-k4-storyboard} 展示了从首轮 GA、统一 F、当前联合求解、交换增强、阈值敏感性到 ALNS 当前最优的完整对照；其中，右下角的节点状态图清楚表明，$K=4$ 最优 ALNS 方案已经把直接确认点数提升到 12 个，仅保留 4 个待人工复核点。图 \ref{fig:problem2-k4-diagnostics} 则从空地完成时间、方案闭环时间、ALNS 种子敏感性与任务切分数量四个角度补充说明：当前最优值 2997.6 s 并非单张路径图上的偶然现象，而是能够在多种指标上同时观察到的结构性改进。

\begin{figure}[H]
\centering
\includegraphics[width=0.95\linewidth]{outputs/figures/c_problem_problem2_alns_merge_visuals.png}
\caption{ALNS 稳定性与跨 $K$ 航线合并启发的图形化对比}
\label{fig:problem2-alns-merge-visuals}
\end{figure}

\begin{figure}[H]
\centering
\includegraphics[width=0.95\linewidth]{outputs/figures/c_problem_problem2_k4_default_vs_alns_paths.png}
\caption{$K=4$ 下默认 F 路径与当前最优 ALNS 路径的对照图}
\label{fig:problem2-k4-default-vs-alns}
\end{figure}

\begin{figure}[H]
\centering
\includegraphics[width=0.95\linewidth]{outputs/figures/c_problem_problem2_k4_air_ground_compare.png}
\caption{$K=4$ 下空中路径与地面 manual\_route 的联合对照图}
\label{fig:problem2-k4-air-ground}
\end{figure}

\begin{figure}[H]
\centering
\includegraphics[width=0.98\linewidth]{outputs/figures/c_problem_problem2_k4_best_solution_storyboard.png}
\caption{$K=4$ 下从默认 F 到 ALNS 当前最优的最优解看板对照}
\label{fig:problem2-k4-storyboard}
\end{figure}

\begin{figure}[H]
\centering
\includegraphics[width=0.98\linewidth]{outputs/figures/c_problem_problem2_k4_best_solution_diagnostics.png}
\caption{$K=4$ 当前最优解的多指标诊断图}
\label{fig:problem2-k4-diagnostics}
\end{figure}

\subsection{全优化流程数据提升历程对比}
为了把当前所有已经跑通的实验结果统一到同一视图下，本文进一步新增一份结构化汇总数据文件 \texttt{c\_problem\_full\_optimization\_report.json}，并同步导出一张适合论文对照的历程表 \texttt{c\_problem\_full\_optimization\_journey\_compare.csv}。后者提炼出 $K=3,4$ 下最具代表性的闭环阶段链，便于直接观察“方法迭代后到底带来了多少收益”。

图 \ref{fig:problem2-journey-flowchart} 则把这条优化链转写为流程图。与单纯堆叠实验表不同，这种表示方式更容易说明当前结果并不是严格单线迭代，而是在“联合求解主线”基础上继续分化出参数敏感性支线和 ALNS 全局搜索支线；其中，最终最优记录来自 $K=4$ 配置下的 ALNS 分支。

\begin{figure}[H]
\centering
\includegraphics[width=0.98\linewidth]{outputs/figures/c_problem_problem2_optimization_journey_flowchart.png}
\caption{问题二优化历程流程图}
\label{fig:problem2-journey-flowchart}
\end{figure}

需要说明的是，问题一输出的主指标是空中阶段最大完工时长，而问题二各轮实验的主指标是总闭环完成时间，因此表 \ref{tab:problem2-full-journey} 只展示问题二中可同口径比较的闭环链路；问题一的原始构造式与层次滚动重构结果则已完整收录到新的汇总 JSON 文件中。

\begin{table}[H]
\centering
\caption{问题二全优化流程数据提升历程对比}
\label{tab:problem2-full-journey}
\small
\setlength{\tabcolsep}{4pt}
\begin{tabularx}{\linewidth}{>{\raggedright\arraybackslash}p{2.5cm}cc>{\raggedright\arraybackslash}X>{\raggedright\arraybackslash}p{2.8cm}}
\toprule
阶段 & $K=3$ 闭环时间/s & $K=4$ 闭环时间/s & 代表新增机制 & 备注 \\
\midrule
首轮 GA 选点闭环 & 4112.2 & 3693.1 & 候选点集全局搜索 & 早期版本 \\
统一 F 引导基线 & 3870.2 & 3540.7 & 状态引导函数 $F(u,v,E)$ & 主线基线 \\
当前联合求解器 & 3799.1 & 3469.5 & 直接确认 + 地面复核统一重算 & 主线更新 \\
交换算子增强 & 3785.5 & 3442.8 & 两条飞行任务间单点交换 & 局部搜索增强 \\
阈值敏感性最优 & 3504.0 & 3349.8 & $\text{multiplier}=0.8$ & 参数支线最优 \\
ALNS 当前最优 & 3752.7 & 2997.6 & destroy-repair 联合搜索 & 历史全量实验最优 \\
ALNS deeper复验 & -- & 3077.4 & seed=37 追加 12 次迭代 & 本轮复验最优 \\
\bottomrule
\end{tabularx}
\normalsize
\end{table}

表 \ref{tab:problem2-full-journey} 呈现出三点值得强调的结论。第一，主线方法本身已经具有稳定增益：从首轮 GA 选点闭环推进到统一 F 引导，再到当前联合求解器，$K=3$ 的总闭环时间累计下降 313.2 s，$K=4$ 累计下降 223.6 s，说明仅靠“更统一的状态评价 + 更一致的空地联动重算”就能得到持续改进。第二，交换算子虽然属于很窄的后处理局部搜索，但在基准阈值不变的前提下仍进一步压缩了 13.6 s 和 26.8 s，证明当前既有航线之间仍存在可被利用的跨任务重配空间。第三，把参数实验、历史最优记录和本轮 deeper 复验都纳入统一视图后，可以更清楚地看到 ALNS 的两层含义：2997.6 s 代表当前已观测到的历史最优闭环解，而 3077.4 s 的 deeper 复验则说明同类 destroy-repair 结构在更长迭代下仍能稳定找到显著优于联合基线的解，只是当前跨 seed 的极值仍有较强波动。

从工程决策角度看，这份历程表已经足以区分“稳健主线”“历史最优”和“复验结果”。若坚持题面原始阈值不变，则 $K=4$ 的联合求解与交换增强已经给出稳定优于早期方案的主线结果；若允许把阈值视为可调参数，则 $\times 0.8$ 的敏感性支线还能继续压缩闭环时间；若追求当前已观测到的最优闭环效率，则历史 ALNS 记录仍是 2997.6 s；而本轮 deeper 复验的 best seed=37 则把总闭环时间稳定压到 3077.4 s，说明扩大搜索预算本身仍然有效，但还需要更多多种子长跑才能把 2997.6 s 这一级别的结果稳定复现。

\section{管理启示与结果预期}
从题目背景出发，模型最终不仅要给出一条最优路径，更要形成可供社区管理者使用的配置建议。预期结果主要包括三类。

第一，随着无人机数量增加，总闭环完成时间将先显著下降，后出现边际收益递减，因此存在一个推荐无人机规模区间。第二，问题二中“是否值得追加悬停”不能只看单点缺口，而应综合考虑当前航线位置、剩余能量和可节省的地面复核代价。第三，当问题二被统一写成 $F(u,v,E)$ 状态引导后，后续 GA、PSO、ACO 的作用就不再是各自独立造一个规则，而是在同一评分架构下学习更优的指引参数。

因此，本文后续实验的目标不是单纯寻找某一组最优参数，而是识别“较优且稳健”的参数区间，为智慧社区巡检系统提供可执行的资源配置方案。

\section{算法实现说明}
为了保证模型可以在竞赛时间内稳定落地，本文已同步搭建 Python 数据脚手架。当前代码模块完成了以下工作：读取 \texttt{2026C数据.xlsx} 中的关键工作表；标准化字段名称；将飞行时间、飞行能耗和地面通行矩阵转化为可直接索引的数值矩阵；并在附件缺少有效地面时间矩阵时，基于物业点坐标、步行速度和绕行系数生成回退矩阵；输出基线统计表，为后续精确模型与启发式算法提供输入。

在此基础上，本文已经实现问题一的可复现实验链，并成功生成不同无人机数量下的结果汇总表与航线明细表。当前实现对应的程序结构包括数据加载模块 \texttt{load\_c\_data.py}、基线统计模块 \texttt{run\_c\_baseline.py}、构造式求解模块 \texttt{run\_problem1\_heuristic.py} 与层次滚动重构模块 \texttt{run\_problem1\_hierarchical.py}。其中，问题一模块能够输出每条航线的访问顺序、每个目标点的悬停时间、每架无人机总作业时间，以及系统最大完工时长、总能耗和负载离散度等指标，并支持在多种构造策略之间做同口径择优。

当前程序已扩展为三层结构：第一层为数据加载与一致性校验，第二层为空中巡检与地面复核的求解模块，第三层为参数实验与图表输出模块。这样的设计便于将论文中的模型表达直接映射到代码实现，降低调试与复现实验的成本。

特别是对问题二而言，本文已经在问题一最优空中路径结果基础上实现了“直接确认判定 + 地面复核路径优化 + 局部交换增强 + ALNS 联合搜索验证”的完整链路。对应的核心脚本包括 \texttt{run\_problem2\_joint.py} 与共享评分模块 \texttt{problem2\_f\_guidance.py}。

在此基础上，围绕问题二与补充实验的脚本分工如下：

\begin{itemize}
	\item \texttt{run\_problem2\_threshold\_sensitivity.py}：批量验证直接确认阈值倍率扰动；
	\item \texttt{run\_problem2\_swap\_experiment.py}：验证跨航线交换算子的收益；
	\item \texttt{run\_problem2\_alns\_joint.py}：负责 destroy-repair 型 ALNS 联合搜索验证；
	\item \texttt{run\_problem2\_joint\_reorder\_f.py}：补做“问题一路径重排 + 问题二 F 引导”的联合实验；
	\item \texttt{run\_problem1\_cross\_k\_merge\_heuristic.py}：验证“高 $K$ 航线反向启发低 $K$ 串行执行”是否能够带来收益。
\end{itemize}

\section{模型优势、局限与后续扩展}
本文初稿框架的优势在于：一是严格对应题面两层任务结构，逻辑递进清晰；二是将空中巡检与地面复核统一到闭环完成时间这一核心指标下，便于做方案比较；三是具有较强工程可实现性，适合结合 AI 辅助快速生成代码、实验和图表。

当前局限首先体现在目标函数的表达仍可进一步细化。现有模型已能刻画总闭环完成时间与直接确认收益之间的基本权衡，但对负载均衡、高优先级点提前完成以及阈值附近的边际收益仍缺乏统一刻画。后续可将多机负载均衡写为时间、能耗与任务点数的归一化离散指标，并引入以完成时刻 $C_s$ 为自变量的优先级收益函数，例如令 $P=\sum_{s\in S} p_s \exp(-\beta C_s)$；同时，对直接确认部分可用平滑收益 $R_s(T_s)=p_s\sigma\!\left(\gamma\left(T_s/h_s-1\right)\right)$ 近似阈值判定，从而使模型不仅比较“是否达到阈值”，还能够区分“接近阈值的追加悬停是否值得”。这样处理后，目标函数可以更精细地描述空中追加停留、人工复核减少和高优先级点优先闭环之间的边际替代关系。

其次，当前求解流程虽然已经显著优于早期两阶段串行方案，但本质上仍以“问题一先形成可行空中路径、问题二再叠加直接确认与人工复核决策”为主，同层联合优化深度仍有提升空间。后续工作的重点应转向“优先级敏感的 F 评分重构 + 更稳定的 ALNS destroy-repair 设计 + 问题一路径重构 + 地面复核路径精化”的同层联合求解，以形成更完整的推荐配置结论。更具体地说，上层不再只决定静态任务分组，而应同时决定任务条数、服务区域与优先级资源倾斜方向；下层则在给定分组内联合优化访问顺序、悬停时间分配、直接确认集合与人工复核路径。

在算法实现上，可将现有 destroy-repair 机制扩展为自适应大邻域搜索，根据“拆分长航线、合并短航线、跨机迁移节点、优先级导向移除和高人工代价点优先回填”等算子的阶段性表现动态更新权重，并把这些算子同时作用于问题一路径结构与问题二闭环决策。这样既能保留小规模实例上的基准最优性，又能提高批量参数实验和更大规模场景下的求解效率与稳定性，也更符合本题“空中巡检资源分配”和“地面复核负担削减”本应协同发生的工程逻辑。

最后，本文现阶段仍建立在静态确定性场景之上，对飞行时间波动、能耗扰动、设备故障和执行过程中的临时事件尚未纳入统一框架。后续可在场景集 $\Omega$ 下引入 $T_{ij}^{u}(\omega)$、$E_{ij}^{u}(\omega)$ 与 $G_{ij}(\omega)$ 的随机扰动，并以期望闭环时间与风险项联合构造鲁棒目标；同时，可结合机会约束限制单次任务能耗超限概率，并在任务执行过程中采用滚动时域重规划，对剩余航线、悬停资源和人工复核顺序进行动态修正。这样能够把当前静态闭环优化进一步推广到鲁棒与动态巡检决策，使最终给出的无人机配置与巡检策略更贴近真实社区运行场景。

\section{结论}
本文已形成一套可以直接支撑竞赛论文写作的 C 题实证链路。问题一方面，本文在补全缺失飞行边的基础上，引入 K-means 起点候选，并比较构造式贪心与层次滚动重构两类算法，最终形成了按无人机数量择优采用的最优问题一方案：当 $K=1,2,3$ 时，构造式方案分别给出 1312.2 s、537.4 s 和 361.3 s 的最大完工时长；当 $K=4$ 时，层次滚动重构进一步将空中阶段压缩到 303.4 s，并把负载标准差降到 14.87 s。问题二方面，本文进一步把“是否值得在某点继续悬停”统一改写为状态引导函数 $F(u,v,E)$，并在此基础上完成了联合求解、交换增强、阈值敏感性和 ALNS 多种子验证，逐步把 $K=4$ 的总闭环时间从 3540.7 s 推进到 2997.6 s。

因此，本文当前阶段已经能够回答题目中的核心工程问题：在现有附件数据下，推荐将无人机数量配置在 $K=3$ 至 $K=4$ 区间；若强调稳健且可复现的主线闭环效率，则 $K=4$ 的联合求解与交换增强已经明显优于早期 GA 与统一 F 基线；若强调当前已观测到的最优闭环效率，则 $K=4$ 配合 ALNS 联合搜索已把总闭环时间进一步压到 2997.6 s，同时将直接确认节点提升到 12 个、待人工复核节点压缩到 4 个，并把任务切分重构为 8 条更短任务。补充实验还表明，阈值倍率 $\times 0.8$ 是当前最有效的参数支线，而简单的跨 $K$ 航线直接合并则不能替代针对目标 $K$ 重新求解。这意味着当前模型已经能较好回答“如何更快完成闭环”，后续工作的重点应转向“优先级敏感的 F 评分重构 + 更稳定的 ALNS destroy-repair 设计 + 问题一路径重构 + 地面复核路径精化”的同层联合求解，以形成更完整的推荐配置结论。

\section{参考文献}
\begin{thebibliography}{99}
\bibitem{r1} 刘卫国，数学建模方法与实践，北京：高等教育出版社，2019，第2版。
\bibitem{r2} Toth P, Vigo D, Vehicle Routing: Problems, Methods, and Applications, Philadelphia: SIAM, 2014, 2nd ed.
\bibitem{r3} Ropke S, Pisinger D, An Adaptive Large Neighborhood Search Heuristic for the Pickup and Delivery Problem with Time Windows, Transportation Science, 40(4): 455--472, 2006.
\bibitem{r4} Dorigo M, Stutzle T, Ant Colony Optimization, Cambridge: MIT Press, 2004, 1st ed.
\bibitem{r5} Kennedy J, Eberhart R, Particle Swarm Optimization, Proceedings of ICNN'95 - International Conference on Neural Networks, 4: 1942--1948, 1995.
\bibitem{r6} Holland J H, Adaptation in Natural and Artificial Systems, Ann Arbor: University of Michigan Press, 1975, 1st ed.
\bibitem{r7} 中国工业和信息化部，民用无人驾驶航空器运行安全管理相关政策文件，\url{https://www.miit.gov.cn/}，访问时间：2026年4月30日。
\end{thebibliography}


\end{document}